\chapter{Architecture}

\section{Vue d'ensemble}

L'architecture suit un modele cloud-native avec separation entre le frontend statique
et le backend API.

% TODO: Remplacer par une capture d'ecran du diagramme d'architecture
% \begin{figure}[H]
%     \centering
%     \includegraphics[width=\textwidth]{architecture-diagram.png}
%     \caption{Diagramme d'architecture globale}
% \end{figure}

\begin{lstlisting}[language={},caption={Architecture simplifiee}]
   Utilisateurs --> CloudFront --> S3 (Angular)
   Utilisateurs --> ALB --> ECS (EC2) --> Express Backend
                                |              |
                           RDS PostgreSQL   DynamoDB (cache)

   SSM Parameter Store <-- secrets
   CloudWatch <-- logs, metriques, alarmes
\end{lstlisting}

\section{Composants}

\subsection{Frontend --- S3 + CloudFront}

L'application Angular est compilee en fichiers statiques et deployee sur un bucket S3.
CloudFront sert de CDN avec~:
\begin{itemize}
    \item Origin Access Control (OAC) pour securiser l'acces au bucket
    \item Redirection HTTPS automatique
    \item Gestion du routage SPA (erreurs 403/404 redirigees vers \texttt{index.html})
\end{itemize}

\subsection{Backend --- ECS (EC2 Launch Type)}

Le backend Express est conteneurise et deploye sur ECS avec le type de lancement EC2.
Ce choix, plutot que Fargate, permet de~:
\begin{itemize}
    \item Configurer les instances EC2 avec Ansible
    \item Demontrer la gestion de l'Auto Scaling Group
    \item Installer des agents de monitoring personnalises
\end{itemize}

Un Application Load Balancer (ALB) distribue le trafic vers les conteneurs
via le port mapping dynamique.

\subsection{Bases de donnees}

\begin{description}
    \item[RDS PostgreSQL] Base relationnelle pour les favoris utilisateur et l'historique
        de recherche. Instance \texttt{db.t3.micro} en single-AZ.
    \item[DynamoDB] Cache des reponses API NASA avec TTL automatique.
        Mode de facturation PAY\_PER\_REQUEST.
\end{description}

\section{Choix d'architecture}

\begin{table}[H]
    \centering
    \begin{tabularx}{\textwidth}{lXl}
        \toprule
        \textbf{Decision} & \textbf{Justification} & \textbf{Alternative} \\
        \midrule
        ECS EC2 vs Fargate & Permet Ansible, plus pedagogique & Fargate \\
        Bridge networking & Standard EC2, port mapping dynamique & awsvpc \\
        Single NAT Gateway & Economie (\~{}32\$/mois) & HA NAT \\
        PriceClass\_100 & Regions US/EU suffisantes & PriceClass\_All \\
        Single-AZ RDS & Projet scolaire, economie & Multi-AZ \\
        \bottomrule
    \end{tabularx}
    \caption{Decisions d'architecture et justifications}
\end{table}

\section{Reseau --- VPC}

\begin{table}[H]
    \centering
    \begin{tabular}{llll}
        \toprule
        \textbf{Sous-reseau} & \textbf{CIDR} & \textbf{AZ} & \textbf{Usage} \\
        \midrule
        public\_a & 10.0.1.0/24 & us-east-1a & ALB, Bastion \\
        public\_b & 10.0.2.0/24 & us-east-1b & ALB (2\textsuperscript{e} AZ) \\
        private\_a & 10.0.10.0/24 & us-east-1a & ECS, RDS \\
        private\_b & 10.0.11.0/24 & us-east-1b & Groupe RDS \\
        \bottomrule
    \end{tabular}
    \caption{Allocation des sous-reseaux}
\end{table}
